\chapter{Hirsutism}

\section*{Definition}
Hirsutism is defined as the presence of terminal (coarse, pigmented) hair in a male-pattern distribution in women, affecting androgen-sensitive areas including the upper lip, chin, chest, back, lower abdomen, and inner thighs. It is distinguished from hypertrichosis, which involves generalized excessive vellus hair growth unrelated to androgen excess, and is quantified using the modified Ferriman-Gallwey score (score of 8 or greater indicates hirsutism).

\section*{Pathophysiology}
Hirsutism arises from increased androgen production, enhanced peripheral conversion of androgens to their active forms, or heightened sensitivity of the pilosebaceous unit to androgens. Terminal hair follicles in androgen-sensitive regions respond to dihydrotestosterone (DHT), derived from testosterone via 5-$\alpha$-reductase, by prolonging the anagen (growth) phase and increasing hair shaft diameter. Ovarian and adrenal sources contribute to circulating androgens, with excess driven by elevated luteinizing hormone, insulin, or ACTH stimulation.

\section*{Etiology}
\begin{itemize}
  \item Ovarian causes: Polycystic ovary syndrome (PCOS, 70--80\% of cases), ovarian hyperthecosis, ovarian androgen-secreting tumors (Sertoli-Leydig cell tumors, granulosa cell tumors)
  \item Adrenal causes: Non-classic congenital adrenal hyperplasia (21-hydroxylase deficiency, 11-beta-hydroxylase deficiency), Cushing syndrome, adrenal carcinoma, adrenal adenoma
  \item Idiopathic hirsutism: Normal circulating androgens with increased skin 5-$\alpha$-reductase activity or androgen receptor sensitivity
  \item Medications: Anabolic steroids, danazol, testosterone, minoxidil, valproate, cyclosporine, corticosteroids (long-term high-dose)
  \item Other endocrine conditions: Hyperprolactinemia, acromegaly, thyroid dysfunction (rare), insulin resistance syndromes
\end{itemize}

\section*{Indications for Evaluation}
Seek care if:
\begin{itemize}
  \item Rapid onset or progression of hirsutism over weeks to months (suggests androgen-secreting tumor)
  \item Hirsutism accompanied by virilization (clitoromegaly, deepening voice, male-pattern baldness, increased muscle mass, breast atrophy)
  \item Onset after menopause
  \item Associated symptoms of Cushing syndrome (central obesity, striae, easy bruising, hypertension)
  \item Irregular menses, infertility, or galactorrhea suggesting PCOS or hyperprolactinemia
\end{itemize}

\section*{Related Symptoms}
\begin{itemize}
  \item Acne
  \item Oligomenorrhea or amenorrhea
  \item Infertility
  \item Androgenetic alopecia
  \item Acanthosis nigricans
  \item Virilization signs
  \item Obesity
  \item Galactorrhea
\end{itemize}
\textbf{Note:} PCOS is the most common cause of hirsutism and is diagnosed by the Rotterdam criteria requiring two of three: oligo-ovulation, clinical or biochemical hyperandrogenism, and polycystic ovaries on ultrasound, with exclusion of other etiologies.

\section*{Self-Care / Home Management}
\begin{itemize}
  \item Cosmetic measures including shaving, waxing, depilatory creams, threading, and bleaching provide immediate but temporary removal.
  \item Topical eflornithine cream (applied twice daily) slows facial hair growth through inhibition of ornithine decarboxylase.
  \item Weight loss in overweight women (BMI $>$ 25) reduces insulin levels and ovarian androgen production, improving hirsutism.
  \item Avoid treatments known to exacerbate hirsutism (valproate, danazol, anabolic steroids).
\end{itemize}

\section*{Diagnostic Approach}
\begin{itemize}
  \item History includes onset, progression, menstruation pattern, fertility, medication use, and family history of hirsutism or PCOS.
  \item Physical examination quantifies hirsutism severity using the modified Ferriman-Gallwey scoring system and assesses for virilization, acne, alopecia, and acanthosis nigricans.
  \item Hormonal evaluation: Total and free testosterone, dehydroepiandrosterone sulfate (DHEA-S), 17-hydroxyprogesterone (screening for non-classic CAH), and prolactin.
  \item Ovarian ultrasound for PCOS diagnosis (polycystic ovary morphology: 12 or more follicles per ovary or volume $>$ 10 cm$^3$).
  \item ACTH stimulation test if 17-hydroxyprogesterone is elevated; overnight dexamethasone suppression test for suspected Cushing syndrome; CT/MRI of adrenals or ovaries if tumor suspected.
\end{itemize}

\section*{Treatment / Management Overview}
\begin{itemize}
  \item Pharmacologic therapy: Combined oral contraceptives (estrogen-progestin) are first-line, suppressing ovarian androgen production and increasing sex hormone-binding globulin.
  \item Anti-androgens: Spironolactone (50--200 mg daily) inhibits androgen receptor and reduces 5-$\alpha$-reductase activity; finasteride blocks 5-$\alpha$-reductase directly.
  \item Insulin sensitizers (metformin) improve ovulatory function and reduce hyperinsulinemia in PCOS but have modest effect on hirsutism alone.
  \item Non-classic CAH is treated with low-dose glucocorticoids (prednisone, dexamethasone) to suppress ACTH-driven adrenal androgen excess.
  \item Permanent hair reduction: Laser hair removal and electrolysis are effective for long-term management of terminal hair in androgen-sensitive areas.
\end{itemize}

\clearpage
